% Banque d'exercices — chapitre 2
% Le chapitre est déterminé par le nom de ce fichier.

\begin{exo}[De la glace froide à l'eau bouillante : ordres de grandeur]{chauffage-glace-eau-ebullition}
\keywords{calorimétrie, changement d'état, fusion, vaporisation, chaleur latente, ordre de grandeur}

\textbf{Données numériques (à la pression atmosphérique) :}
\[
\begin{aligned}
m &= 1{,}0\,\text{kg}, & T_i &= -100\,{}^\circ\text{C},\\
c_{\text{glace}} &= 2{,}1\times10^3\,\text{J}\!\cdot\!\text{kg}^{-1}\!\cdot\!\text{K}^{-1},
& L_f &= 3{,}34\times10^5\,\text{J}\!\cdot\!\text{kg}^{-1},\\
c_{\text{eau}} &= 4{,}18\times10^3\,\text{J}\!\cdot\!\text{kg}^{-1}\!\cdot\!\text{K}^{-1},
& L_v &= 2{,}26\times10^6\,\text{J}\!\cdot\!\text{kg}^{-1}.
\end{aligned}
\]
À la pression atmosphérique, on chauffe progressivement un kilogramme de glace,
initialement à $-100\,{}^\circ\text{C}$.

\begin{enumerate}
  \item Calculer l'énergie nécessaire pour amener la glace de $-100\,{}^\circ\text{C}$ à $0\,{}^\circ\text{C}$.
  \item Calculer l'énergie nécessaire pour faire fondre entièrement la glace à $0\,{}^\circ\text{C}$.
  \item Calculer l'énergie nécessaire pour chauffer l'eau liquide de $0\,{}^\circ\text{C}$ à $100\,{}^\circ\text{C}$.
  \item Calculer l'énergie supplémentaire nécessaire pour vaporiser entièrement l'eau à $100\,{}^\circ\text{C}$.
  \item Quelle étape demande le plus d'énergie ? Comparer l'énergie du chauffage de $0\,{}^\circ\text{C}$ à $100\,{}^\circ\text{C}$ d'un kilogramme d'eau à l'énergie potentielle d'une masse $m$ tombant de dix mètres de haut. Quelle masse $m$ faut-il ainsi faire tomber pour égaler cette énergie ?
  \item On considère maintenant le processus inverse. Une masse $m_v$ de vapeur d'eau se condense sur un pare-brise, sans que l'eau liquide obtenue se refroidisse ensuite. Exprimer la chaleur $Q_{\text{lib}}$ libérée lors de la condensation, puis la calculer pour $m_v=1{,}0\times10^{-2}\,\text{kg}$.
\end{enumerate}

\begin{indication}
Pour chauffer une phase sans changement d'état, utiliser $Q=mc\Delta T$. Pour un changement d'état à température constante, utiliser $Q=mL$. L'énergie potentielle de pesanteur d'une masse $m$ située à une hauteur $h$ vaut $E_p=mgh$ ; on prendra $g=9{,}81\,\text{m}\!\cdot\!\text{s}^{-2}$.
\end{indication}

\begin{solution}
\textbf{1.} Le chauffage de la glace jusqu'à sa température de fusion demande
\[
Q_1=mc_{\text{glace}}\Delta T
=1{,}0\times2{,}1\times10^3\times100
=2{,}1\times10^5\,\text{J}.
\]

\textbf{2.} Pendant la fusion, la température reste égale à $0\,{}^\circ\text{C}$ :
\[
Q_2=mL_f=1{,}0\times3{,}34\times10^5
=3{,}34\times10^5\,\text{J}.
\]

\textbf{3.} Pour chauffer l'eau liquide de $0$ à $100\,{}^\circ\text{C}$,
\[
Q_3=mc_{\text{eau}}\Delta T
=1{,}0\times4{,}18\times10^3\times100
=4{,}18\times10^5\,\text{J}.
\]

\textbf{4.} La vaporisation complète demande encore
\[
Q_4=mL_v=1{,}0\times2{,}26\times10^6
=2{,}26\times10^6\,\text{J}.
\]
Cette étape demande donc à elle seule une énergie de l'ordre de quelques mégajoules.

\textbf{5.} La vaporisation est de loin l'étape la plus coûteuse :
\[
Q_4=2{,}26\times10^6\,\text{J}
>Q_3=4{,}18\times10^5\,\text{J}
>Q_2=3{,}34\times10^5\,\text{J}
>Q_1=2{,}1\times10^5\,\text{J}.
\]
La vaporisation exige notamment environ $(2{,}26\times10^6)/(4{,}18\times10^5)\approx5{,}4$ fois plus d'énergie que le chauffage de l'eau liquide de $0$ à $100\,{}^\circ\text{C}$.

Pour une masse $m$ tombant d'une hauteur $h=10\,\text{m}$, l'énergie potentielle libérée vaut
\[
E_p=mgh.
\]
On cherche la masse telle que $mgh=Q_3$. Ainsi,
\[
m=\frac{Q_3}{gh}
=\frac{4{,}18\times10^5}{9{,}81\times10}
\approx4{,}26\times10^3\,\text{kg}.
\]
Il faudrait donc faire tomber d'une hauteur de dix mètres une masse d'environ $4{,}3\times10^3\,\text{kg}$ pour libérer autant d'énergie que celle nécessaire au chauffage d'un kilogramme d'eau de $0$ à $100\,{}^\circ\text{C}$. Avec l'approximation $g\approx10\,\text{m}\!\cdot\!\text{s}^{-2}$, on trouve directement $m\approx4{,}2\times10^3\,\text{kg}$.

\textbf{6.} La condensation est le processus inverse de la vaporisation : la vapeur cède au pare-brise la chaleur latente
\[
Q_{\text{lib}}=m_vL_v.
\]
Pour $m_v=1{,}0\times10^{-2}\,\text{kg}$,
\[
Q_{\text{lib}}=1{,}0\times10^{-2}\times2{,}26\times10^6
=2{,}26\times10^4\,\text{J}.
\]
Ainsi, la condensation d'une faible masse de vapeur peut libérer une quantité de chaleur non négligeable. Cette chaleur est reçue par le pare-brise et son environnement.
\end{solution}

\end{exo}

\begin{exo}[Un degré de trop en Méditerranée]{odg-rechauffement-mediterranee}
\keywords{ordre de grandeur, Méditerranée, température de surface, capacité thermique, pluie, cyclone, chaleur latente}

On observe lors de certains épisodes chauds un excès de température de surface de la mer Méditerranée de l'ordre de $1\,{}^\circ\text{C}$. Une température mesurée exactement à la surface ne suffit pas à déterminer une énergie : il faut connaître la profondeur sur laquelle l'anomalie s'étend. On adopte ici un modèle simple dans lequel les dix premiers mètres de la mer sont uniformément plus chauds de $1\,{}^\circ\text{C}$.

\textbf{Données numériques :}
\[
\begin{aligned}
A_{\text{Méd}}&=2{,}5\times10^{12}\,\text{m}^2,
& h&=10\,\text{m},
& \Delta T&=1\,\text{K},\\
\rho_{\text{eau}}&=1{,}0\times10^3\,\text{kg}\!\cdot\!\text{m}^{-3},
& c_{\text{eau}}&=4{,}0\times10^3\,\text{J}\!\cdot\!\text{kg}^{-1}\!\cdot\!\text{K}^{-1},
& L_v&=2{,}5\times10^6\,\text{J}\!\cdot\!\text{kg}^{-1}.
\end{aligned}
\]

Pour fixer les idées, on modélise un épisode de pluie méditerranéen très intense par une hauteur d'eau $H=1{,}0\times10^{-1}\,\text{m}$ tombant sur une surface $S=1{,}0\times10^{10}\,\text{m}^2$. On admet également qu'un cyclone tropical mature libère, par condensation de la vapeur d'eau, une puissance thermique de l'ordre de $P_{\text{cyc}}=6\times10^{14}\,\text{W}$.

\begin{enumerate}
  \item Estimer la masse d'eau contenue dans la couche superficielle d'épaisseur $h=10\,\text{m}$.
  \item Estimer l'énergie thermique $E_{\text{Méd}}$ correspondant à un excès uniforme de température de $1\,{}^\circ\text{C}$ dans cette couche. En donner l'ordre de grandeur en joules.
  \item Calculer la masse d'eau qui tombe lors de l'épisode pluvieux décrit ci-dessus, puis l'énergie thermique $E_{\text{pluie}}$ libérée dans l'atmosphère lors de la condensation de cette eau. Combien de tels épisodes représentent, en énergie, l'excès thermique calculé à la question précédente ?
  \item Pendant quelle durée un cyclone tropical mature doit-il libérer de la chaleur à la puissance $P_{\text{cyc}}$ pour atteindre cette énergie ? Donner d'abord le résultat en secondes, puis en jours pour l'interpréter.
  \item Si toute l'énergie excédentaire de la couche superficielle était évacuée uniformément pendant $\tau=2{,}6\times10^6\,\text{s}$, calculer la puissance thermique moyenne correspondante, puis le flux thermique moyen par unité de surface de la Méditerranée.
  \item Pourquoi ces comparaisons énergétiques ne permettent-elles pas d'affirmer que l'excès thermique de la Méditerranée sera intégralement converti en pluie ou en énergie cyclonique ?
\end{enumerate}

\begin{indication}
Le volume de la couche superficielle vaut $Ah$ et sa masse $\rho Ah$. L'énergie libérée par la condensation d'une masse $m$ d'eau vaut $mL_v$. Enfin, une puissance est une énergie divisée par une durée, et un flux thermique une puissance divisée par une surface.
\end{indication}

\begin{solution}
\textbf{1.} Le volume des dix premiers mètres est
\[
V=A_{\text{Méd}}h
=2{,}5\times10^{12}\times10
=2{,}5\times10^{13}\,\text{m}^3,
\]
et leur masse vaut
\[
M=\rho_{\text{eau}}V
=10^3\times2{,}5\times10^{13}
=2{,}5\times10^{16}\,\text{kg}.
\]

\textbf{2.} L'énergie thermique excédentaire est
\[
E_{\text{Méd}}=Mc_{\text{eau}}\Delta T
=2{,}5\times10^{16}\times4{,}0\times10^3\times1
=1{,}0\times10^{20}\,\text{J}.
\]
L'ordre de grandeur recherché est donc $10^{20}\,\text{J}$. Le résultat est proportionnel à la profondeur choisie : une anomalie limitée au premier mètre représenterait dix fois moins d'énergie, tandis qu'une anomalie atteignant $50\,\text{m}$ en représenterait cinq fois plus.

\textbf{3.} Le volume puis la masse de pluie sont
\[
V_{\text{pluie}}=SH=1{,}0\times10^{10}\times1{,}0\times10^{-1}=1{,}0\times10^9\,\text{m}^3,
\qquad
m_{\text{pluie}}=\rho_{\text{eau}}V_{\text{pluie}}=10^{12}\,\text{kg}.
\]
La condensation de cette eau libère approximativement
\[
E_{\text{pluie}}=m_{\text{pluie}}L_v
=10^{12}\times2{,}5\times10^6
=2{,}5\times10^{18}\,\text{J}.
\]
Le rapport vaut
\[
\frac{E_{\text{Méd}}}{E_{\text{pluie}}}
=\frac{10^{20}}{2{,}5\times10^{18}}
\approx40.
\]
L'excès thermique modélisé représente donc le même ordre de grandeur que la chaleur latente libérée par quelques dizaines de tels épisodes de pluie très intense.

\textbf{4.} La durée recherchée est
\[
\tau_{\text{cyc}}=\frac{E_{\text{Méd}}}{P_{\text{cyc}}}
=\frac{1{,}0\times10^{20}}{6\times10^{14}}
\approx1{,}7\times10^5\,\text{s}
\approx1{,}9\,\text{jours}.
\]
L'énergie excédentaire estimée est ainsi comparable à la chaleur latente libérée en environ deux journées par un cyclone tropical mature. Il s'agit de chaleur dégagée par la condensation, et non de la seule énergie cinétique des vents, qui est beaucoup plus faible.

\textbf{5.} La puissance moyenne serait
\[
P=\frac{E_{\text{Méd}}}{\tau}
=\frac{10^{20}}{2{,}6\times10^6}
\approx3{,}9\times10^{13}\,\text{W}.
\]
Répartie sur toute la surface de la Méditerranée, elle correspondrait au flux
\[
\Phi=\frac{P}{A_{\text{Méd}}}
=\frac{3{,}9\times10^{13}}{2{,}5\times10^{12}}
\approx16\,\text{W}\!\cdot\!\text{m}^{-2}.
\]

\textbf{6.} Ce calcul compare des réservoirs et des transferts d'énergie, mais ne décrit pas leur dynamique. Une partie de la chaleur de la mer est transportée par les courants, mélangée vers des couches plus profondes ou rayonnée vers l'espace ; seule une fraction est transférée à l'atmosphère par évaporation, conduction et rayonnement. La formation d'une pluie extrême ou d'un cyclone dépend aussi de l'humidité, de la stabilité de l'atmosphère, des vents et de leur cisaillement, de la pression et de la circulation à grande échelle. Une énergie disponible importante peut donc favoriser certains phénomènes sans être intégralement ni automatiquement convertie en un événement météorologique donné.
\end{solution}

\end{exo}

\begin{exo}[Mélange de deux masses d'eau]{calorimetrie-melange-eau}
\keywords{calorimétrie, mélange, capacité thermique, chaleur spécifique, Joseph Black, température d'équilibre}

On verse $m_1 = 0{,}200\,\text{kg}$ d'eau à $T_1 = 80\,^\circ\text{C}$ dans un récipient contenant déjà $m_2 = 0{,}300\,\text{kg}$ d'eau à $T_2 = 15\,^\circ\text{C}$. Le récipient est un calorimètre idéal~: on néglige toute perte de chaleur vers l'extérieur ainsi que la capacité thermique du récipient lui-même. On rappelle la relation
\[
Q = mc\Delta T,
\]
où $c = 4{,}18\times10^3\,\text{J}\!\cdot\!\text{kg}^{-1}\!\cdot\!\text{K}^{-1}$ est la capacité thermique massique de l'eau liquide.

\begin{enumerate}
  \item Justifier que la chaleur cédée par l'eau chaude est intégralement reçue par l'eau froide~: écrire cette conservation sous la forme d'une équation faisant intervenir $m_1$, $m_2$, $c$, $T_1$, $T_2$ et la température d'équilibre $T_{eq}$.
  \item En déduire l'expression littérale de $T_{eq}$, puis sa valeur numérique.
  \item Calculer la quantité de chaleur $Q$ cédée par l'eau chaude au cours de ce mélange.
  \item De tels mélanges d'une même substance permettent-ils de déterminer la valeur de $c$ ? Justifier.
\end{enumerate}

\begin{indication}
Le calorimètre étant isolé et rigide, l'énergie interne totale $U_1 + U_2$ se conserve : la chaleur cédée par l'un est reçue par l'autre au signe près.
\end{indication}

\begin{solution}
\textbf{1.} Le calorimètre étant isolé (parois adiabatiques et rigides vers l'extérieur), l'énergie interne totale du système $\{$eau chaude, eau froide$\}$ est conservée~: ce que perd l'une, l'autre le gagne exactement. Avec la convention $Q = mc\Delta T$ appliquée à chaque masse d'eau entre son état initial et l'équilibre,
\[
m_1 c (T_{eq} - T_1) + m_2 c (T_{eq} - T_2) = 0.
\]

\textbf{2.} Le coefficient $c$ se simplifie (même substance des deux côtés), et on obtient
\[
T_{eq} = \frac{m_1 T_1 + m_2 T_2}{m_1 + m_2}
= \frac{0{,}200 \times 80 + 0{,}300 \times 15}{0{,}200 + 0{,}300}
= 41\,^\circ\text{C}.
\]

\textbf{3.} La chaleur cédée par l'eau chaude en refroidissant de $80\,^\circ\text{C}$ à $41\,^\circ\text{C}$ vaut
\[
Q = m_1 c (T_1 - T_{eq})
=0{,}200\times4{,}18\times10^3\times(80-41)
\approx3{,}26\times10^4\,\text{J}.
\]
On vérifie que l'eau froide reçoit exactement cette même quantité en s'échauffant de $15\,^\circ\text{C}$ à $41\,^\circ\text{C}$~:
\[
m_2c(T_{eq}-T_2)
=0{,}300\times4{,}18\times10^3\times26
\approx3{,}26\times10^4\,\text{J}.
\]

\textbf{4.} Non. Pour deux masses de la même substance, le même coefficient $c$ multiplie les deux termes du bilan
\[
m_1c(T_{eq}-T_1)+m_2c(T_{eq}-T_2)=0
\]
et se simplifie. La température d'équilibre mesurée ne dépend donc pas de $c$ : elle est uniquement la moyenne des températures initiales pondérée par les masses. Un tel mélange ne permet pas de déterminer $c$. Pour mesurer une capacité thermique massique par la méthode des mélanges, il faut introduire une substance de capacité thermique connue, comme dans l'exercice suivant.
\end{solution}

\end{exo}

\begin{exo}[Détermination de la capacité thermique d'un métal par la méthode des mélanges]{calorimetrie-methode-melanges}
\keywords{calorimétrie, méthode des mélanges, capacité thermique massique, calorimètre, chaleur spécifique, valeur en eau}

On souhaite mesurer, comme le faisait déjà Joseph Black au XVIII\textsuperscript{e}~siècle, la capacité thermique massique $c_m$ d'un métal inconnu par la \emph{méthode des mélanges}~: on chauffe un échantillon de ce métal, de masse $m = 0{,}150\,\text{kg}$, jusqu'à $T_m = 95\,^\circ\text{C}$, puis on le plonge rapidement dans un calorimètre contenant $m_e = 0{,}250\,\text{kg}$ d'eau ($c_e = 4{,}18\times10^3\,\text{J}\!\cdot\!\text{kg}^{-1}\!\cdot\!\text{K}^{-1}$) initialement à $T_e = 18\,^\circ\text{C}$. On mesure, une fois l'équilibre thermique atteint, $T_{eq} = 22\,^\circ\text{C}$. Dans un premier temps, on néglige la capacité thermique propre du récipient du calorimètre (parois, agitateur, thermomètre).

\begin{enumerate}
  \item Écrire, comme à l'exercice précédent, la conservation de l'énergie du système isolé $\{$métal, eau$\}$, en faisant apparaître $c_m$ comme seule inconnue.
  \item En déduire l'expression littérale de $c_m$, puis sa valeur numérique. À quel métal usuel cette valeur correspond-elle approximativement (comparer aux valeurs de référence~: aluminium $\approx 9{,}0\times10^2\,\text{J}\!\cdot\!\text{kg}^{-1}\!\cdot\!\text{K}^{-1}$, fer $\approx 4{,}5\times10^2$, cuivre $\approx 3{,}9\times10^2$, plomb $\approx 1{,}3\times10^2$) ?
  \item En réalité, le récipient du calorimètre lui-même absorbe une partie de la chaleur cédée par le métal. On modélise cet effet par une \emph{valeur en eau} $\mu = 1{,}5\times10^{-2}\,\text{kg}$~: tout se passe comme si le calorimètre se comportait comme une masse supplémentaire $\mu$ d'eau. Reprendre le calcul de $c_m$ en tenant compte de $\mu$, et commenter le sens de la correction (le $c_m$ corrigé est-il plus grand ou plus petit que celui de la question 2~?).
  \item Pourquoi est-il essentiel, expérimentalement, de plonger l'échantillon \emph{rapidement} dans le calorimètre, et de bien l'isoler thermiquement de l'air ambiant pendant la mesure ?
\end{enumerate}

\begin{indication}
Le métal cède la chaleur $Q_m = mc_m(T_m - T_{eq})$, intégralement reçue par l'eau (et, question 3, par le calorimètre assimilé à une masse d'eau supplémentaire $\mu$)~: $Q_m = (m_e+\mu)c_e(T_{eq}-T_e)$.
\end{indication}

\begin{solution}
\textbf{1.} Le système $\{$métal, eau$\}$ étant isolé, la chaleur cédée par le métal en refroidissant de $T_m$ à $T_{eq}$ est intégralement reçue par l'eau qui s'échauffe de $T_e$ à $T_{eq}$~:
\[
mc_m(T_m - T_{eq}) = m_e c_e (T_{eq} - T_e).
\]

\textbf{2.} On isole $c_m$~:
\[
c_m = \frac{m_e c_e (T_{eq}-T_e)}{m(T_m-T_{eq})}
= \frac{0{,}250\times4{,}18\times10^3\times(22-18)}{0{,}150\times(95-22)}
\approx3{,}82\times10^2\,\text{J}\!\cdot\!\text{kg}^{-1}\!\cdot\!\text{K}^{-1}.
\]
Cette valeur est très proche de celle du cuivre ($\approx3{,}9\times10^2\,\text{J}\!\cdot\!\text{kg}^{-1}\!\cdot\!\text{K}^{-1}$).

\textbf{3.} En tenant compte de la valeur en eau du calorimètre, l'eau \emph{et} le récipient reçoivent ensemble la chaleur cédée par le métal~:
\[
c_m = \frac{(m_e+\mu) c_e (T_{eq}-T_e)}{m(T_m-T_{eq})}
= \frac{(0{,}250+0{,}015)\times4{,}18\times10^3\times4}{0{,}150\times73}
\approx4{,}05\times10^2\,\text{J}\!\cdot\!\text{kg}^{-1}\!\cdot\!\text{K}^{-1}.
\]
La correction augmente légèrement $c_m$~: en négligeant la valeur en eau du calorimètre à la question 2, on sous-estimait la quantité de chaleur réellement cédée par le métal (une partie servait aussi à chauffer le récipient), et donc on sous-estimait $c_m$.

\textbf{4.} Toute la méthode repose sur l'hypothèse d'un système isolé pendant la durée de la mesure~: plonger l'échantillon rapidement limite le temps pendant lequel il peut échanger de la chaleur avec l'air ambiant avant d'atteindre l'eau, et bien isoler le calorimètre limite les pertes vers l'extérieur pendant que l'équilibre s'établit. Toute fuite thermique non comptabilisée fausserait le bilan de la question 1 et donc la valeur déduite de $c_m$ -- exactement le genre de soin expérimental qu'apportait déjà Black, puis Lavoisier et Laplace avec leur calorimètre à glace.
\end{solution}

\end{exo}

\begin{exo}[Calories alimentaires et puissance métabolique]{calorimetrie-calories-alimentaires}
\keywords{calorie, kilocalorie, Calorie alimentaire, équivalent mécanique, puissance métabolique, unités}

L'étiquette nutritionnelle indique qu'un adulte consomme en moyenne environ $2000\,\text{kcal}$ par jour. La kilocalorie est une unité non-SI encore utilisée en nutrition ; on donne directement sa conversion vers le SI : $1\,\text{kcal}=4{,}18\times10^3\,\text{J}$.

\begin{enumerate}
  \item Convertir cette ration quotidienne en joules.
  \item En déduire la puissance métabolique moyenne de cet adulte, en watts, en supposant cette énergie consommée uniformément pendant $\Delta t=8{,}64\times10^4\,\text{s}$.
  \item Une barre énergétique affiche « $210\,\text{kcal}$ ». Quelle masse d'eau initialement à $20\,^\circ\text{C}$ pourrait-on porter à ébullition ($100\,^\circ\text{C}$) avec cette même quantité d'énergie, si elle était intégralement convertie en chaleur ? On prendra $c_{\text{eau}} = 4{,}18\times10^3\,\text{J}\!\cdot\!\text{kg}^{-1}\!\cdot\!\text{K}^{-1}$.
  \item Le corps humain n'est bien sûr pas un poêle~: il ne restitue pas directement cette énergie sous forme de chaleur pure. Qu'est-ce que ce calcul mesure malgré tout correctement, et qu'est-ce qu'il ne prend pas en compte ?
\end{enumerate}

\begin{indication}
Pour la question 2, utiliser $\mathcal P=E/\Delta t$. Pour la question 3, convertir d'abord l'énergie indiquée sur l'étiquette en joules, puis isoler $m$ dans $Q=mc\Delta T$.
\end{indication}

\begin{solution}
\textbf{1.} En unités SI,
\[
E=2000\times4{,}18\times10^3
=8{,}36\times10^6\,\text{J}.
\]

\textbf{2.} La puissance moyenne vaut
\[
\mathcal{P} = \frac{8{,}36\times10^6}{8{,}64\times10^4} \approx 97\,\text{W}.
\]
Un être humain au repos consomme donc, en moyenne, à peu près autant de puissance qu'une ampoule à incandescence -- un ordre de grandeur qui surprend souvent, et qui illustre bien l'intérêt de convertir les calories alimentaires en unités physiques usuelles.

\textbf{3.} L'énergie de la barre, convertie en joules, vaut
\[
Q=210\times4{,}18\times10^3
\approx8{,}78\times10^5\,\text{J}.
\]
Avec $\Delta T=80\,\text{K}$,
\[
m = \frac{Q}{c\Delta T}
= \frac{8{,}78\times10^5}{4{,}18\times10^3\times80}
\approx2{,}63\,\text{kg}.
\]
Une seule barre énergétique représente donc, en ordre de grandeur, de quoi porter environ $2{,}5\,\text{kg}$ d'eau du robinet à ébullition.

\textbf{4.} Ce calcul mesure correctement le contenu \emph{énergétique total} de l'aliment, tel que mesuré par calorimétrie directe (combustion complète de l'aliment dans un calorimètre, à la manière de Lavoisier et Laplace) -- c'est très exactement la définition physique de la kilocalorie alimentaire. Il ne prend en revanche pas en compte le fait que le métabolisme ne convertit pas cette énergie en chaleur pure~: une partie sert à un travail chimique et mécanique utile (contraction musculaire, synthèse de molécules, maintien du potentiel électrique des cellules, etc.), une partie n'est de toute façon pas absorbée par l'organisme (digestion incomplète), et l'énergie effectivement libérée sous forme de chaleur ne l'est pas de façon uniforme dans le temps comme le suppose l'hypothèse de la question 2.
\end{solution}

\end{exo}




\begin{exo}[Mayer et Joule : l'équivalent mécanique de la chaleur]{mayer-equivalent-mecanique}
\keywords{Mayer, Joule, premier principe, équivalent mécanique, chaleur, travail, histoire}

Une chute d'eau de hauteur $h=50\,\text{m}$ transporte un débit massique $\dot m=1{,}0\times10^3\,\text{kg}\!\cdot\!\text{s}^{-1}$. On admet, comme Mayer (1842) et Joule après lui, que toute l'énergie mécanique perdue par l'eau au cours de sa chute se retrouve intégralement sous forme de chaleur -- l'idée étant qu'\emph{aucune} énergie ne disparaît, mais qu'elle change seulement de forme, de mécanique à thermique.

\begin{enumerate}
\item Exprimer puis calculer la puissance mécanique $\mathcal{P}$ perdue par ce flux d'eau, en fonction de $\dot m$, $h$ et $g=9{,}81\,\text{m}\!\cdot\!\text{s}^{-2}$.
\item On isole maintenant, par la pensée, $1\,\text{kg}$ de cette eau qui chute de la hauteur $h$ puis s'immobilise, et on suppose que la totalité de son énergie mécanique perdue reste emmagasinée dans cette même eau sous forme de chaleur (elle ne se dilue pas dans un bassin extérieur). Exprimer, en joules, l'énergie $W$ ainsi convertie.
\item Par définition historique, une calorie (cal) est la quantité de chaleur nécessaire pour élever de $1\,\text{K}$ la température de $1{,}0\times10^{-3}\,\text{kg}$ d'eau liquide. On mesure expérimentalement que ce même kilogramme d'eau voit sa température s'élever de $\Delta T\approx0{,}117\,\text{K}$ après la chute. Exprimer, \emph{en calories, sans utiliser le joule ni aucune constante exprimée en $\text{J}\!\cdot\!\text{kg}^{-1}\!\cdot\!\text{K}^{-1}$}, la quantité de chaleur $Q_{\text{cal}}$ correspondante.
\item En déduire l'équivalent mécanique de la chaleur $J = W/Q_{\text{cal}}$, exprimé en $\text{J}\!\cdot\!\text{cal}^{-1}$, et comparer à la valeur moderne ($4{,}184\,\text{J}\!\cdot\!\text{cal}^{-1}$).
\item Comment ce résultat préfigure-t-il le premier principe ?
\end{enumerate}

\begin{indication}
Question 1 : $\mathcal{P}=\dot mgh$. Question 3 : déterminer combien de masses élémentaires de $1{,}0\times10^{-3}\,\text{kg}$ sont contenues dans $1\,\text{kg}$. La définition historique de la calorie ne fait intervenir aucun joule -- c'est précisément ce qui rend la question 4 non circulaire.
\end{indication}

\begin{solution}
\textbf{Question 1.}
\[
\mathcal{P}=\dot mgh
=1{,}0\times10^3\times9{,}81\times50
\approx4{,}9\times10^5\,\text{W}.
\]
C'est l'ordre de grandeur de la puissance d'une petite centrale électrique.

\textbf{Question 2.} Pour $m=1\,\text{kg}$,
\[
W=mgh=1\times9{,}81\times50=4{,}905\times10^2\,\text{J}.
\]

\textbf{Question 3.} La masse de $1\,\text{kg}$ contient
\[
\frac{1\,\text{kg}}{1{,}0\times10^{-3}\,\text{kg}}=1000
\]
fois la masse de référence intervenant dans la définition de la calorie. Par conséquent,
\[
Q_{\text{cal}}=1000\times0{,}117\approx117\,\text{cal}.
\]

\textbf{Question 4.} $J = W/Q_{\text{cal}} = 490{,}5 / 117 \approx 4{,}19\,\text{J}\!\cdot\!\text{cal}^{-1}$, en excellent accord avec la valeur moderne $4{,}184\,\text{J}\!\cdot\!\text{cal}^{-1}$. Le point essentiel est que ce calcul ne présuppose \emph{à aucun moment} la valeur numérique de $J$ : le travail $W$ est mesuré en unités purement mécaniques ($\text{kg}$, $\text{m}$, $\text{s}$), et la chaleur $Q_{\text{cal}}$ en unités purement thermiques, définies par un protocole expérimental sur l'eau (calories) sans référence au travail. C'est cette indépendance des deux mesures qui rend la détermination de $J$ non triviale et physiquement significative -- au contraire d'un calcul qui utiliserait d'emblée la capacité thermique de l'eau exprimée en $\text{J}\!\cdot\!\text{kg}^{-1}\!\cdot\!\text{K}^{-1}$, ce qui reviendrait à supposer déjà connu le résultat que l'on cherche à établir.

\textbf{Question 5.} Mayer, puis Joule par des expériences bien plus systématiques (la roue à palettes brassant de l'eau, 1845--1850), établissent ainsi que travail mécanique et chaleur sont deux formes d'échange d'une seule et même grandeur conservée : l'énergie. La constante $J$ n'a de sens que tant que l'on mesure la chaleur et le travail dans des unités historiquement différentes ; une fois cette équivalence acceptée, on peut choisir une unité commune (le joule) pour les deux, et $J$ disparaît purement et simplement de l'écriture moderne du premier principe, $\Delta U = W + Q$.
\end{solution}
\end{exo}
